\begin{solapartat}
Disminució de la massa:

\punts{0,25}
\[ \Delta m = 4\,m\left({}^{1}_{1}p\right) - m\left({}^{4}_{2}He\right) = 4,765\times 10^{-29}\ \mathrm{kg} \]
Energia alliberada per nucli d'heli:

\punts{0,25}
\[ E_{at,Heli} = \Delta m\, c^2 = 4,765\times 10^{-29}\times(3,00\times 10^{8})^2 = 4,29\times 10^{-12}\ \mathrm{J} \]

\punts{0,5} Per formar un nucli d'heli calen 2 molècules d'aigua; per tant, el nombre total de nuclis d'heli que es poden formar és:
\[ n = N_A\,\frac{n_{mols\ H2O}}{2} = 5,12\times 10^{24} \]

\punts{0,25} I, finalment, l'energia total que es pot extreure del got d'aigua és:
\[ E_{Tot} = n\,\Delta m\, c^2 = 5,12\times 10^{24}\times 4,29\times 10^{-12} = 2,20\times 10^{13}\ \mathrm{J} \]
\end{solapartat}

\begin{solapartat}
\punts{0,2} Una partícula $\alpha$ és un nucli d'heli (és un àtom d'heli totalment ionitzat): ${}^{4}_{2}He$ o ${}^{4}_{2}\alpha$.

\punts{0,7}
\[ {}^{6}_{3}Li + {}^{1}_{0}n \rightarrow {}^{4}_{2}He + {}^{3}_{1}H \]

\punts{0,15} A partir de l'energia que es pot extreure de l'aigua mitjançant la cadena protó-protó i el consum energètic per persona, podem comprovar que l'energia continguda a l'aigua és superior al consum d'una persona durant mil anys.

No obstant això, en els reactors de fusió nuclear no es produeix la cadena protó-protó, sinó que l'energia s'obté de la fusió del triti amb el deuteri. La quantitat de triti natural és ínfima i s'ha d'obtenir artificialment a partir de la fissió del liti. Així doncs, no és cert que amb un simple got d'aigua es pugui produir l'energia que consumeix una família de quatre persones durant tota la vida. Per tant, no es pot considerar una font d'energia tan il·limitada com la quantitat d'aigua del nostre planeta.

\punts{0,2} Per una altra banda, la fusió del deuteri i del triti genera neutrons d'elevada energia que poden convertir en radioactius els materials del reactor. Certament, la fusió genera menys residus nuclears que la fissió, però no es pot considerar totalment neta.
\end{solapartat}
